\section{Programación Entera}

\subsection{Introducción}

Un problema de \textbf{Programación Lineal Entera} (PLE) es un problema de optimización de una función lineal en presencia
de restricciones lineales con la condición añadida de que alguna o todas las variables deben ser enteras. Cuando todas las
variables están obligadas a ser enteras entonces se habla de un problema de \textbf{Programación Entera Pura} y si solo son
algunas será un problema de \textbf{Programación Entera Mixta}. En ocasiones las variables enteras están limitadas a tomar
valores 0 o 1, a estas variables las denominaremos binarias. Si esta condición es extensiva a todas las variables entonces
hablamos de un problema de \textbf{Programación 0-1 Pura}, si solo afecta a alguna de las variables presentes, en este caso
hablamos de un problema de \textbf{Programación 0-1 Mixta}.

\vspace*{2mm}

Los problemas de Programación Entera en general presentan grandes dificultades computacionales, pues el método Simplex deja
de funcionar. Por ello se recurre a otros métodos de resolución, como los algoritmos \textit{branch-and-bound} (BNB) o 
algoritmos de planos de corte.

\vspace*{2mm}

Este capítulo está dedicado a la formulación de problemas de Programación Lineal con presencia de variables enteras o 
binarias y, en especial, a la modelización de condiciones o restricciones lógicas. Estas últimas restricciones son capaces
de convertir un problema continuo en otro entero, por ello se suele considerar que el $90\%$ de los problemas reales de 
programación terminan siendo de Programación Entera Pura o Mixta.


\begin{flushright}
    \textit{A partir de ahora se usarán ejemplos para explicar los distintos tipos de problemas a estudiar}
\end{flushright}

\subsection{Programación Entera Pura}

\begin{ejemplo}
    Una empresa está planificando la inversión a realizar en los próximos 3 trimestres en 5 proyectos de investigación. En 
    la siguiente tabla aparece detallado el presupuesto disponible por trimestre, la inversión que se debe realizar por
    proyecto y trimestre, y el beneficio que se espera obtener con cada proyecto. Todos los datos aparecen en millones de 
    euros.

    \begin{center}
        \begin{tabular}{|c|ccc|c|}
        \hline
        \multirow{2}{*}{\textbf{Proyecto}} & \multicolumn{3}{c|}{\textbf{Trimestre}} & \multirow{2}{*}{\textbf{Beneficio}} \\
         & 1 & 2 & 3 & \\
        \hline
        1 & 5 & 1 & 8 & 20 \\
        2 & 4 & 7 & 10 & 40 \\
        3 & 7 & 9 & 2 & 20 \\
        4 & 3 & 4 & 1 & 15 \\
        5 & 8 & 6 & 10 & 30 \\
        \hline
        \shortstack{\textbf{Presupuesto} \\ \textbf{disponible}} & 25 & 25 & 25 & \multicolumn{1}{c}{} \\
        \cline{1-4}
        \end{tabular}
    \end{center}

    Se desea conocer los proyectos en los que se debe invertir para maximizar el beneficio obtenido.
\end{ejemplo}

Obsérvese que no buscamos la cantidad a invertir en cada proyecto, sino que la única decisión que debemos tomar es si se
invierte en un proyecto o no con el fin de maximizar el beneficio sin superar el presupuesto disponible. Se consideran 
entonces las variables 0-1 siguientes

\[
    x_i=
    \begin{cases}
        1 & \text{si se invierte en el proyecto } i \\
        0 & \text{si no se invierte en el proyecto } i
    \end{cases}
\]

\newpage

La formulación matemática del problema es la siguiente

\begin{alignat}{4}
    \text{maximizar:}   & \quad & 20x_1+40x_2+20x_3+15x_4+30x_5     & \quad &           \\
    \text{sujeto a:}    & \quad & 5x_1+4x_2+3x_3+7x_4+8x_5\leq 25   &       &           \\
                        & \quad & x_1+7x_2+9x_3+4x_4+6x_5\leq 25    &       &           \\
                        & \quad & 8x_1+10x_2+2x_3+x_4+10x_5\leq 25  &       &           \\
                        & \quad & x_i\in \{0, 1\}                   &       & \forall i=1,\dots,5          
\end{alignat}

Todas las variables presentes en este problema están restringidas a tomar los valores enteros 0 o 1, luego estamos ante un
problema de Programación 0-1 Pura. La solución es invertir en todos los proyectos salvo el 5, generando un beneficio de
95 millones de euros.

\vspace*{2mm}

\noindent Imaginemos que la empresa impone algunas de las siguientes condiciones sobre la inversión

\begin{enumerate}
    \item Se debe invertir, como mucho, en tres proyectos.
    \item Se debe invertir en el proyecto 1 o en el 5 (o en ambos, or lógico).
    \item Si se invierte en el proyecto 4 se debe invertir en el 3.
    \item Si se invierte en los proyectos 2 o 3 se debe invertir en el proyecto 5.
\end{enumerate}

Formalizar la primera restricción es sencillo, como toman valores 0-1, contar el número de inversiones es equivalente a sumar
las variables, por lo que la restricción queda tan sencilla como $x_1+\cdots+x_5\leq 3$. La segunda tiene truco, pues si bien
podemos obligar la presencia de una variable con una restricción del tipo $x_i\geq 1$, si incorporamos a la formulación las 
restricciones $x_1\geq 1$, $x_5\geq 1$ estaríamos obligando a la presencia de ambas conjunta pero sin permitir que aparezcan
de forma individual. Una representación posible y correcta sería $x_1+x_5\geq 1$. La tercera condición hace que una variable
dependa de la otra, y esto lo debemos reflejar en la restricción de forma que si $x_4$ toma el valor 1, $x_3$ también por lo
que la restricción queda $x_3\geq x_4$ (se puede ver que no hay problema para que cuando $x_4=0$, $x_3$ sea 1). Por último,
la cuarta condición no puede ser representada de manera análoga a la segunda, pues algo del tipo $x_5\geq x_2+x_3$ cuando 
ambas variables sean 1 $x_5$ debería, como mínimo, valer 2, por lo que rompe la condición de variable binaria. Es por esto
que resulta mucho más sencillo formalizarla como dos restricciones $x_5\geq x_2$ y $x_5\geq x_3$.

\subsection{Programación Entera Mixta}

\begin{ejemplo}
    Una pequeña empresa puede producir 4 tipos de artículos y estima que para el próximo periodo, podrá vender todo lo que 
    fabrica de los productos 1 y 3, pero como mucho tendrá demanda para 30 unidades del producto 2 y 10 del producto 4. La 
    empresa tiene capacidad para producir un máximo de 110 artículos de cualquier tipo y combinación. A continuación se 
    detalla el coste de producción por unidad producida, el coste fijo de producción y el precio de venta unitario de 
    cada artículo

    \begin{center}
        \begin{tabular}{|c|c|c|c|}
        \hline
        \textbf{Producto} & \textbf{Coste de} & \textbf{Coste fijo} & \textbf{P.V.P.} \\
         & \textbf{producción} & & \\
        \hline
        1 & 300 & 60 & 450 \\
        2 & 450 & 20 & 650 \\
        3 & 500 & 70 & 600 \\
        4 & 200 & 60 & 450 \\
        \hline
        \end{tabular}
    \end{center}
    
    Formular el modelo matemático para determinar las unidades que se deben producir de cada artículo para maximizar los
    beneficios de la empresa.
\end{ejemplo}

Si consideramos únicamente las variables 

\[
    x_i=\text{unidades del producto } i\text{ que debe fabricar el próximo periodo}
\]

la función de beneficios no podrá formularse correctamente. Ténganse en cuenta que los costes fijos no dependen de la cantidad
sino de la decisión de producir o no, por lo que se necesitan más variables que hagan referencia a este aspecto

\[
    y_i=
    \begin{cases}
        1 & \text{si se fabrica el artículo } i \\
        0 & \text{si no se fabrica el artículo } i
    \end{cases}
\]

Por lo tanto se formula el problema de la siguiente forma

\begin{alignat}{5}
    \text{maximizar:}   & \quad & 150x_1+200x_2+100x_3+250x_4-60y_1-20y_2-70y_3-60y_4   & \quad &           \\
    \text{sujeto a:}    & \quad & x_1+x_2+x_3+x_4\leq 110                               &       &           \\
                        & \quad & x_2\leq 30                                            &       &           \\
                        & \quad & x_4\leq 10                                            &       &           \\
                        & \quad & x_i\geq 0                                             &       &           \\
                        & \quad & y_i\in \{0, 1\}                                       &       & \forall i=1,\dots,4          
\end{alignat}

Ahora bien, está clara que la solución no respeta la relación entre las variables de producción y de decisión, pues ``podemos''
(las restricciones lo permiten) fabricar productos sin pagar los costes fijos. Esto es debido a que no se han incluido
relaciones entre las variables. 

\newpage

Sea $M$ una constante arbitrariamente grande. Entonces la restricción $x_i\leq My_i$ asegura que cuando $y_i=1$ entonces 
$x_i>0$, o reciprocamente cuando $y_i=0$ entonces $x_i\leq 0$. La idea básica es que al asegurar que $x_i$ sea 0 cuando se
toma la decisión de no comprar, entonces $x_i$ será mayor que 0 cuando se tome la decisión de comprar (suena un poco absurdo
así explicado pero mi intención era, en efecto, reducirlo al absurdo). No tenemos que preocuparnos de que $x_i$ sea más grande
 que $M$ en la solución óptima pues $M$ es lo suficientemente grande.

\vspace*{2mm}

\noindent Supongamos ahora que el fabricante impone además las siguientes condiciones

\begin{enumerate}
    \item Si se produce del artículo 1 se debe producir del artículo 2
    \item Si se produce del artículo 4 no debe producirse nada del artículo 3.
    \item Se debe producir obligatoriamente de los artículos 2, 3 y 4.
    \item Si se producen más de 30 unidades del artículo 1, se deben producir por lo menos 10 unidades del artículo 2.
    \item Si se producen más de 30 unidades del artículo 1, se deben producir como mucho 10 unidades del artículo 2.
\end{enumerate}

La primera condición se puede expresar como $x_1 > 0 \implies x_2 > 0$ que se traduce como $x_1\leq My_1$ y $x_2\geq my_1$,
cumpliendose así que $x_1>0\implies y_1=1 \implies x_2\geq m$. La segunda se resuelve simplemente haciendo que ambas variables
de decisión no puedan estar presentes al mismo tiempo mediante la restricción $y_3+y_4\leq 1$. En cuanto a la tercera es tan
simple como añadir 3 restricciones que obliguen a elegir esos artículos ($y_2=y_3=y_4=1$) o lo que es equivalente 
$y_2+y_3+y_4=3$. Las siguientes 2 restricciones implican crear una nueva variable $z$ definida como $z=0$ si $x_1\leq 30$ y
$z=1$ si $x_1>30$ y añadiendo la restricción $x_1-30\leq Mz$ con $M$ suficientemente grande, lo que permite definir la 
cuarta condición mediante la restricción $x_2\geq 10z$ y la quinta mediante la restricción $x_1-10\leq M(1-z)$ para $M$
suficientemente grande.

\vspace*{5mm}

\noindent Generalizaremos alguna de las ideas expresadas en los ejemplos anteriores.

\subsubsection{Restricciones ``si entonces'' en programación mixta}

En general, la condición $f(x_1,\dots,x_n)>0 \implies g(x_1,\dots,x_n)\geq k$ se formula del modo siguiente. 
Definimos la variable lógica $z$ tal que

\[
    z=
    \begin{cases}
        1 & \text{si } f(x_1,\dots,x_n)>0 \\
        0 & \text{si } f(x_1,\dots,x_n)\leq0
    \end{cases}
\]

\noindent e incorporaremos al modelo las condiciones

\[
    \begin{array}{l}
        f(x_1,\dots,x_n)\leq Mz \\
        g(x_1,\dots,x_n)\geq kz \\
        z \in {0, 1}
    \end{array}
\]

\newpage

Reciprocamente, la condición $f(x_1,\dots,x_n)>0 \implies g(x_1,\dots,x_n)\leq k$ se formula incorporando al modelo 
las condiciones

\[
    \begin{array}{l}
        f(x_1,\dots,x_n)\leq Mz \\
        g(x_1,\dots,x_n)-k\geq M(1-z) \\
        z \in {0, 1}
    \end{array}
\]

\subsection{Problemas de programación entera típicos}

Se habla de los distintos tipos de problemas de programación entera en las siguientes secciones de los
apuntes, incuyendo la sección de \textbf{Extras}, la cual contiene problemas más específicos.

